\documentclass[tikz,border=6pt]{standalone}
\usepackage{ctex}
\usepackage{amsmath}
\usetikzlibrary{calc, arrows.meta, decorations.pathreplacing}

\begin{document}
\begin{tikzpicture}[scale=1.15, thick]

% 抛物线焦点弦示意图
% 抛物线 y^2 = 4x，焦点 F(1,0)，准线 x = -1
% 焦点弦 AB 过 F，A(x1,y1), B(x2,y2)
%
% 取斜率 k=1 的焦点弦：直线 y=x-1，代入 y^2=4x
% (x-1)^2=4x => x^2-6x+1=0 => x=(6±sqrt(32))/2=3±2√2
% A: x1=3+2√2≈5.828, y1=x1-1=2+2√2≈4.828  (超出画布)
% 改取斜率 k=2：直线 y=2(x-1)=2x-2，代入 (2x-2)^2=4x
% 4x^2-8x+4=4x => 4x^2-12x+4=0 => x^2-3x+1=0
% x=(3±√5)/2 => x1=(3+√5)/2≈2.618, y1=2x1-2≈3.236
%             x2=(3-√5)/2≈0.382, y2=2x2-2≈-1.236

\def\xA{2.618}
\def\yA{3.236}
\def\xB{0.382}
\def\yB{-1.236}

% 坐标轴
\draw[->, gray, thin] (-1.6, 0) -- (5.5, 0) node[right, font=\footnotesize] {$x$};
\draw[->, gray, thin] (0, -2.8) -- (0, 4.5) node[above, font=\footnotesize] {$y$};
\node[font=\footnotesize, below right] at (0,0) {$O$};

% 准线 x = -1
\draw[gray, dashed, thin] (-1, -2.6) -- (-1, 4.2);
\node[font=\footnotesize, below] at (-1, -2.6) {准线 $x=-p$};

% 抛物线 y^2=4x
\draw[black, thick, domain=0:4.2, samples=80, smooth]
  plot ({(\x)^2/4}, \x);
\draw[black, thick, domain=-2.8:0, samples=60, smooth]
  plot ({(\x)^2/4}, \x);

% 焦点 F(1,0)
\coordinate (F) at (1, 0);
\fill[red] (F) circle [radius=2pt];
\node[font=\footnotesize, below right] at (F) {$F(p,0)$};

% 焦点弦 AB（延长至画布边界）
\draw[blue, thick] (-0.1, -2.2) -- (3.2, 4.4);

% 交点 A, B
\coordinate (A) at (\xA, \yA);
\coordinate (B) at (\xB, \yB);
\fill[black] (A) circle [radius=2pt];
\fill[black] (B) circle [radius=2pt];
\node[font=\footnotesize, right] at (A) {$A(x_1,y_1)$};
\node[font=\footnotesize, right] at (B) {$B(x_2,y_2)$};

% 准线垂足 A' 和 B'（A, B 到准线的垂足）
\coordinate (Ap) at (-1, \yA);
\coordinate (Bp) at (-1, \yB);
\fill[gray!60] (Ap) circle [radius=1.5pt];
\fill[gray!60] (Bp) circle [radius=1.5pt];
\node[font=\footnotesize, left] at (Ap) {$A'$};
\node[font=\footnotesize, left] at (Bp) {$B'$};

% 焦半径（A 到 F，B 到 F）
\draw[red!70, thick] (A) -- (F);
\draw[red!70, thick] (B) -- (F);

% A 到准线的水平线
\draw[gray, dashed, thin] (A) -- (Ap);
\draw[gray, dashed, thin] (B) -- (Bp);

% |AF| 标注
\node[red, font=\footnotesize] at ({(\xA+1)/2+0.2}, {\yA/2+0.6}) {$|AF|$};
% |BF| 标注
\node[red, font=\footnotesize] at ({(\xB+1)/2+0.1}, {\yB/2-0.35}) {$|BF|$};

% 韦达定理框
\node[draw, rounded corners, fill=yellow!15, font=\small, align=left,
      text width=5.0cm] at (4.5, 1.8)
  {%
   \textbf{韦达定理（$k=2$，$y^2=4x$）}\\[3pt]
   直线 $y=2(x-1)$ 代入\\
   $\Rightarrow x^2-3x+1=0$\\[2pt]
   $x_1+x_2=3$\\
   $x_1 x_2 = 1$\\[3pt]
   $y_1+y_2=2(x_1+x_2)-4=2$\\
   $y_1 y_2=(2x_1-2)(2x_2-2)$\\
   $\quad=4(x_1 x_2 - x_1 - x_2 + 1)$\\
   $\quad=4(1-3+1)=-4$%
  };

% 调和关系框
\node[draw, rounded corners, fill=blue!8, font=\small, align=left,
      text width=5.2cm] at (4.5, -1.4)
  {%
   \textbf{调和关系（焦点弦特技）}\\[3pt]
   $|AF|=x_1+p,\quad |BF|=x_2+p$\\[2pt]
   $\dfrac{1}{|AF|}+\dfrac{1}{|BF|}$
   $=\dfrac{x_1+x_2+2p}{(x_1+p)(x_2+p)}$\\[4pt]
   由韦达：$x_1 x_2=\dfrac{p^2}{1}$（$y^2=2px$）\\[2pt]
   化简得：$\boxed{\dfrac{1}{|AF|}+\dfrac{1}{|BF|}=\dfrac{2}{p}}$%
  };

% 顶点
\fill[gray] (0,0) circle [radius=1.8pt];
\node[font=\footnotesize, below left] at (0,0) {};

\end{tikzpicture}
\end{document}
